\section{zkSnark Design}
\subsection{Overview}
This section introduces 
\foldpot,
an efficient folding scheme
that is tailored for efficiently verifying
the discharging algorithms introduced earlier.
\foldpot, as its name suggests, mixes
the design from several
recent folding schemes, including
SuperNova \cite{SuperNova}, \footnote{SuperNova uses a linear
folding model as Nova \cite{Nova}. 
Here we assume heuristically
that there is no known attack over the number of steps being folded.
In \cite[Remark 6.3]{BCCT13}, it was pointed out that stronger assumption on
extractor size/time extractor (e.g., additive over prover) leads
to polynomial step bound (instead of the constant bound if
assuming polynomial extractor for a general SNARK scheme for each step
in a recursive SNARK composition scheme).
This was shown in recent workk \cite{AGMRecursion}.
It gives a contrived example
which shows that when blackbox knowledge extractor run time becomes
superpolynomial after $\log(\lambda)$ steps, soundness can be broken.
However, it also shows that under a modified AGM model (where prover keys
are uniformly distributed), Nova/SuperNova is secure given that
Pedersen vector commitment is used. Here
we use the sequential folding scheme for the problem size
we deal with, which incurs maximum $2^{17}$ steps, though
PCD tree based folding/accuulation is possible with further
development work.}
Cyclefold \cite{CycleFold}, Mangrove \cite{Mangrove} as well as
finger-printing schemes such as Mirage+ \cite{AIPRam}, and the
idea of NARK compiler from ProtoStar \cite{ProtoStar}.
We also assilate the idea of finger-printing techniques,
which were practiced in many previous works \cite{Mirage,AIPRam,Zkreg},
and adapt it in the folding schemes.

We first introduce the design goals that lead to the design
of \foldpot.  First, we would like to discharge a 
large collection of files (e.g., all Linux
executables). It is known that folding/accumulation 
provides an overall more efficient solution
than proof agregation schemes SNARKPack \cite{SNARKPack} 
and GIPA \cite{Bunz21} both asympototically (at verifier) and
concretely (at prover). Given that discharging of each individual file
needs to stack discharging steps non-uniformly,
we decide to adopt SuperNova \cite{SuperNova} for its 
support of non-uniform sub-circuits. It is appealing to provide
concretely small and constant size proof, especially verifiable
on-chain (e.g., ETH), CycleFold \cite{CycleFold} provides
an efficient decider proof generation over half pairing groups
BN254 + Grumpkin, so that $O(1)$ final decider proof can be
generated over BN254.
Finally, the system relies heavily on finger-printing techniques,
and our circuit model adapts the one from Mirage+ \cite{AIPRam}.
We then integrate the idea of commit-and-fold from 
Mangrove \cite{Mangrove}: commitments of partial segment of
witnesses can be added into committed R1CS instance, and the check
can be delayed to final decider proof. Leveraging on Mangrove's
commit-and-fold scheme, we develop low-cost techniques for
passing data between non-uniform circuits.

Our contributions are: (1) A compressing scheme for interactive R1CS
model \cite{AIPRam} that results in less communication rounds (and hence
less recursion overhead); By instantiating the commit-and-fold
scheme in Mangrove \cite{Mangrove}, we are able to produce: 
(2) An efficient input/output passing scheme
between non-uniform step-circuits, which has considerable benefits
over general RAM techniques for streaming circuits; (3) 
A constant-size batching-proof protocol that handles the batch
proof problem for non-uniform problem statements. 

\subsection{Circuit Model}
To integrate the finger-printing technique in folding schemes,
we mix the I-R1CS model for Algebraic Interactive Protocol (AIP)
in \cite{AIPRam} and the commited R1CS relation of Nova/SuperNova
\cite{Nova}. The 
finger-printing technique divides witness into multiple segments, where
subsequent segments depend on the Fiat-Shamir public coin 
randoms that are generated from the hash of previous segments.
This technique has been widely deployed in real-world proving
systems such as Halo2 \cite{Halo2}.
The challenge here in folding
is that {\em the number of segments determines the number of
commitments}, which linearly increases the cost of the prover
for final decider. In our context,
one additional commitments costs at least
$5$ million R1CS per non-uniform circuit.
We therefore seek to reduce the number of segments,
and present a $\Sigma$-IR1CS model (which is a weaker model
than the general I-R1CS) but allows only one segment commitment.

We first motivate the model. Consider the $\cp$ algorithm, which
given an input word $s$, returns a collection of signatures
that cannot be discharged by $\cp$. The algorithm walks $s$
through ACDFA, from the acceptance path, extracts
a multi-set of acceptance states, compressing it into
a unique set, performing a join relation with a table that
maps from final states to signatures, and then converting 
the multi-set into a standard set. The algorithm can be
regarded as a sequential composition of a constant number
of component interactive protcols that heavily rely on
public-coin challenges from the verifier. For instance,
to prove two multi-sets $\myvec{a} \in \F^n$ and $\myvec{b} \in \F^n$ 
are permutation of each other, the verifier samples a challenge
$\beta$ and verifies, which is widely used, e.g., in lookup arguments
such as plookup \cite{plookup}:\footnote{CHECK TODO: see if this
ref is accurate}
\[
	\prod_{i=1}^n (\myvec{a}_i + \beta)  =
	\prod_{i=1}^n (\myvec{b}_i + \beta) 
\]

Such sequential composition of interactive protocols is
nicely captured by the notion of AIP and I-R1CS \cite{AIPRam},
where an I-RICS instance has a constant number of message
exchanges, and a final R1CS check (which is similar to the
NARK verifier in ProtoStar \cite{ProtoStar}, with the
restriction of degree $2$). 
In \cite{AIPRam}, the multi-round interaction is handled
by a multi-segment variation of Groth16 \cite{Groth16}
and commit-and-prove extension of Mirage \cite{Mirage}, because
verifier can perform the Fiat-Shamir to generate public coin
challenges out-of-circuit.

In our context, the circuit has to be folded. Generating a commitment
for each round of messages to be exchanged is expensive, because
it increases recursion overhead. Instead, we use a 
restricted I-R1CS (called $\Sigma$-I-R1CS) that consists of
only two messages from the prover to verifier, one at
the beginning, consisting of the concatenation of problem statements
and all messages that are independent of public coins, and
the last message consists of all intermediate messages that
depend on random challenges. We show that a restricted AIP bound by
a certain ``join" property can always be
compiled into a 3-move I-R1CS, and then using the 
ProtoStar compiler \cite{ProtoStar}, compiled into a NARK.
We motivate our construction using an example.

\begin{example} \label{examp:ir1cs}
Assume there is a lookup table $\myvec{T} \subseteq \F$ which defines
a range that excludes $0$ (e.g., a set of final states of an ACDFA,
which ranges from $1$ to $2^{k}-1$ for some $k$).
Given $\myvec{a} \in \myvec{T}^{m}$ and $\myvec{b} \in \myvec{T}^n$,
consider the task of proving that $\myvec{a}$ is the {\em set support}
of multi-set $\myvec{b}$, that is: $\myvec{a}$ includes all the elements of
$\myvec{b}$, and $\myvec{b}$ includes all of $\myvec{a}$, and there is no duplicate element in $\myvec{a}$, assuming that all elements of
$\myvec{a}$ and $\myvec{b}$ belong to $\myvec{T}$.

We can denote it as a relation 
$(\myvec{a}, \myvec{b}, \myvec{T}; \myvec{w}) \in R$, where
a four element tuple $(m,n,2^{k}-1, 0)$ defines $R$'s structure (length
of each message vector).
Here $\myvec{a}$, $\myvec{b}$, $\myvec{T}$ are the
problem statement that 
is visible by both the prover and verifier.
$\myvec{w}$ (separated using \texttt{;} from the statement)
is the secret witness of prover. For this example,
$\myvec{w} = \emptyset$.
We call $R$ an {\em algebraic relation with a fixed structure} 
as all elements involved
in the relation are field elements.\footnote{We note that 
group elements can always be represented non-natively given a 
target field by bit-splitting.} 
We sometimes write $(\myvec{a}, \myvec{b}, \myvec{T}) \in R$
for short when context is clear.

Note that relation $R$ has an interesting property: it is a 
{\em natural join} of three 
relations: $R_1$, $R_1'$, and $R_2$. Here, 
$R_1(\myvec{x},\myvec{y};\myvec{w}_1)$ asserts a lookup relation
that all elements of $\myvec{x}$ belongs to $\myvec{y}$.
$R_1'$ is exactly the same as $R_1$ except changing
the config by swapping the size of the first 
two problem statement elements.
$R_2(\myvec{x}, \myvec{T};\myvec{w}_2)$ asserts that
all elements of $\myvec{x}$ are unique (assuming all 
$\myvec{x}_i \in \myvec{T}$).
Apparently $(\myvec{a}, \myvec{b}, \myvec{T}) \in R$ if and only if 
there exists $\myvec{w}_1$, $\myvec{w}_2$, and $\myvec{w}_3$ s.t.
$\myvec(\myvec{a}, \myvec{b}; \myvec{w}_1) \in R_1$,
$\myvec(\myvec{b}, \myvec{a}; \myvec{w}_1) \in R_1$,
and $\myvec(\myvec{a}, \myvec{T}; \myvec{w}_2) \in R_2$.\footnote{In 
real implementation, we allow $\myvec{a}$ and $\myvec{b}$ to
contain $0$ as dummy filler elements to make the circuits
more general purpose at slightly more constriant cost.}

For $R$, each of its sub-relation has a constant move special
sound protocol (similar to those defined in ProtoStar \cite{ProtoStar}).
For instance, Consider $R_1$, it has a $3$-move special sound protocol
by \cite{Hab22}, as shown in $\proto_1$ in Figure \ref{fig:2shots}. 
Similarly, let $\proto_1'$ be the protocol for $R_1'$.
Its basic idea
is that if the lookup relation holds, the sum of logarithmic derivative
of the elements of two vectors will match. Protocol
$\proto_2$ for $R_2$ is essentially running $\proto_1$ to prove
that the difference for each pair of elements of $\myvec{x}$ is
a positive integer, excluding $0$ (because $0 \not\in \myvec{T}$),
thus proving $\myvec{x}$ is sorted in ascending order, hence
a set without duplicates. One can see that each corresponds to
an I-R1CS protocol where a degree $2$ function is performed
at the end for making checks. It is known that degree $2$ functions
can be expressed in R1CS.

Sequential composition of these protocols is costly in folding.
Instead, we can compile them into $\proto$ as shown in Figure\ \ref{fig:2shots}.
Basically, it combines all the first message from all sub-protocols
in its first message, keeping the structure of public coin verifier
challenges, and then combine all ``response" messages from sub-protocols
in one last message, and then projecting each individual messages
from the combined messages, and run the degree $2$ check functions
correspondingly.  We call this compiled protocol {\em $\Sigma$-I-R1CS} 
because it has 3 moves.

We can then apply the comiler technique in ProtoStar
\cite[Section 3.2]{ProtoStar}, to compile a
$\Sigma$-I-R1CS into a committed protocol
by sending commitments of the first and
third message, and add additional checks to open them (as shown in
$\proto_{\textrm{cm}}$ in Figure\ \ref{fig:2shots}).
One might wonder why sending the commitments and then
immediately open them up. The reason is that these commitments
are included as part of the committed instances
to be folded, and the opening check only performed at the final decider proof,
{\em only once} for many instances that are folded.
\end{example}

Let $\rohash$  denotes the hash function to instantiate 
the random-oracle.
We summarize our discussion and the main finding below.

\begin{definition} An algebraic relation $R$ is an NP relation
over $\F$ with a fixed structure.
We call $R(\myvec{x})$ {\em decomposible by joins} of a constant number of
relations $R_1, ... R_k$ if all of the following
are satisfied:
(1) Each $R_i$ has a special-sound
$\Sigma$-protocol, and
(2) Denote $\myvec{x}$ as $\myvec{x}^{(k+1)}$.
There exists a join function $\varphi$ s.t.
$\varphi$ is a DNF (disjunctive normal form) of 
equality constraints of form $\myvec{x}^{(i)}_u = \myvec{x}^{(j)}_v$,
where $0\leq i,j \leq k+1$.
For any $(\myvec{x}, \myvec{w})$:
$(\myvec{x}; \myvec{w}) \in R$ if and only if
there exists 
$(\myvec{x}^{(1)}; \myvec{w}^{(1)}) \in R_1$,
$(\myvec{x}^{(2)}; \myvec{w}^{(2)}) \in R_2$,
...
$(\myvec{x}^{(k)}; \myvec{w}^{(k)}) \in R_k$, 
and $\varphi(\myvec{x}^{(1)}, ..., \myvec{x}^{(k)},\myvec{x}) = \texttt{true}$.
\end{definition}

We note that there can be more expressive form of expressing the
join relation $\varphi$, but this form sufficies for our applications.
Note that $\varphi$ is degree $1$, hence it can be encoded in R1CS.

\begin{figure*}
%\resizebox{4.5in}{4.5in}{
\resizebox{0.90\textwidth}{!}{
	\begin{tikzpicture}
 \tikzstyle{every node}=[font=\fontsize{20}{20}\selectfont]
\matrix (m)[matrix of nodes, column sep=0.5mm,row  sep=8mm, nodes={draw=none, anchor=center,text depth=0pt} ]{
 %row 1 
 & $\proto_1$ for $R_1 = \{(\myvec{x}\in \F^{m},\myvec{y}\in \F^{n}) ~|~\forall a\in \myvec{x}: a\in \myvec{y} \}$ & \\[-4mm]
 %row 2 
 Prover & & Verifier\\[-4mm]
 %row 3
 Compute $\myvec{m}^{(1)} \in \F^n$ s.t. $\myvec{m}^{(1)}_i$ & & \\ [-7mm]
 %row 4
 is the count of $\myvec{y}_i$ in $\myvec{x}$ & &  \\ [-7mm]
 %row 5
 & $\myvec{x}$, $\myvec{y}$, $\myvec{m}^{(1)}$ & \\  [-7mm]
 %row 6
 & & $r_1 \pick \F$ \\ [-7mm]
 %row 7
 & $r_1$ & \\ [-7mm]
 %row 8
 Compute $\myvec{h}^{(1)} = \left(\frac{1}{\myvec{x}_i+r_1}\right)_{i=1}^m$ & & \\ [-7mm]
 %row 9
 Compute $\myvec{g}^{(1)} = \left(\frac{\myvec{m}^{(1)}_i}{\myvec{y}_i+r_1}\right)_{i=1}^n$ & & \\ [-7mm]
 %row 10 
 & $\myvec{h}^{(1)}$, $\myvec{g}^{(1)}$ & \\ [-7mm]
 %row 11 
 & & Compute $1/0 \leftarrow f_1(\myvec{x}, \myvec{y}, \myvec{h}^{(1)}, \myvec{g}^{(1)})$: \\ [-7mm]
 %row 12 
 &  & $\forall_i\in [m]: \myvec{h}^{(1)}_i(\myvec{x}+r_1) = 1$ \\ [-7mm]
 %row 13 
 & & $\forall_i\in [n]: \myvec{g}^{(1)}_i(\myvec{y}+r_1) = \myvec{m}^{(1)}_i$ \\ [-7mm]
 %row 14 
 & & $\sum_{i=1}^m \myvec{h}^{(1)}_i = \sum_{i=1}^n \myvec{g}^{(1)}_i$  \\[7mm]
 %row 15 
 & $\proto$ for $R_2 = \{(\myvec{x}\in \F^{m},\myvec{T}=\left(i\right)_{i=1}^{2^k-1}) ~|~\forall i\in [m-1]: \myvec{x}_{i+1}>\myvec{x}_{i} \}$ & \\[-4mm]
 %row 16 
 Prover & & Verifier\\[-4mm]
 %row 17 
 Compute $\myvec{m}^{(2)} \in \F^{2^{k}-1}$ s.t. $\myvec{m}^{(2)}_i$ & & \\ [-7mm]
 %row 18 
 is the count of $\myvec{T}_i$ in $\big(\myvec{x}_{i+1}-\myvec{x}_{i}\big)_{i=1}^{m-1}$ & &  \\ [-7mm]
 %row 19 
 & $\myvec{x}, \myvec{T}, \myvec{m}^{(2)}$ & \\ [-7mm]
 %row 20 
 & & $r_2 \pick \F$ \\ [-7mm]
 %row 21 
 & $r_2$ & \\ [-7mm]
 %row 22 
 Compute $\myvec{h}^{(2)} = \left(\frac{1}{\myvec{x}_{i+1}-\myvec{x}_{i}+r_2}\right)_{i=1}^{m-1}$ & & \\ [-7mm]
 %row 23 
 Compute $\myvec{g}^{(2)} = \left(\frac{\myvec{m}^{(2)}_i}{\myvec{T}_i+r_2}\right)_{i=1}^{2^k-1}$ & & \\ [-7mm]
 %row 24 
 & $\myvec{h}^{(2)}$, $\myvec{g}^{(2)}$ & \\ [-7mm]
 %row 25 
 & & Compute $1/0 \leftarrow f_2(\myvec{x}, \myvec{T}, \myvec{h}^{(2)}, \myvec{g}^{(2)})$: \\ [-7mm]
 %row 26 
 &  & $\forall_i\in [m-1]: \myvec{h}^{(2)}_i (\myvec{x}_{i+1} - \myvec{x}_{i} +r_2) = 1$ \\ [-7mm]
 %row 27 
 & & $\forall_i\in [2^k-1]: \myvec{g}^{(2)}_i (\myvec{T}_{i}+r_2) = \myvec{m}^{(2)}_i$ \\ [-7mm]
 %row 28 
 & & $\sum_{i=1}^{m-1} \myvec{h}^{(2)}_i = \sum_{i=1}^{2^k-1} \myvec{g}^{(2)}_i$  \\[7mm]
 %row 29 
 & $\proto$ for $R = \{(\myvec{x}\in \F^{m},\myvec{y}\in \F^{n},\myvec{T}=\left(i\right)_{i=1}^{2^k-1}) ~|~\forall a\in \myvec{x}: a\in \myvec{y} \And \forall b\in \myvec{y}: b\in \myvec{x} \And \forall i\neq j: \myvec{x}_i \neq \myvec{x}_j\}$ & \\[-4mm]
 %row 30 
 Prover & & Verifier\\[-4mm]
 %row 31 
 Compute $\myvec{m}^{(1)}, \myvec{m}^{(1)'}, \myvec{m}^{(2)}$ as $\proto_1$, $\proto_1'$, and $\proto_2$ & & \\ [-7mm]
 %row 32
 & $\myvec{x}, \myvec{y}, \myvec{m}^{(1)}, \myvec{m}^{(1)'}, \myvec{T}, \myvec{m}^{(2)}$ & \\[-7mm]
 %row 33
 & & $r_1, r_1', r_2 \pick \F$ \\[-7mm]
 %row 34
 & $r_1$, $r_1'$, $r_2$ & \\[-7mm]
 %row 35
 Compute 
 $\myvec{h}^{(1)}, \myvec{g}^{(1)},$ 
 $\myvec{h}^{(1)'}, \myvec{g}^{(1)'},$ 
 $\myvec{h}^{(2)}, \myvec{g}^{(2)}$ & & \\[-7mm]
 %row 36
 in $\proto_1$, $\proto_1'$, and $\proto_2$. & & \\[-7mm]
 %row 37
 & 
 	$\myvec{h}^{(1)}, \myvec{g}^{(1)},$
 	$\myvec{h}^{(1)'}, \myvec{g}^{(1)'},$ 
	$\myvec{h}^{(2)}, \myvec{g}^{(2)}$ & \\[-7mm]
 %row 38
 & & Check $f_1$, $f_1'$, $f_2$  \\
 %row 39  ----- Committed Protocol
 & $\proto_{\textrm{cm}}$ for $R = \{(\myvec{x},\myvec{y},\myvec{T}) ~|~\forall a\in \myvec{x}: a\in \myvec{y} \And \forall i\neq j: \myvec{x}_i \neq \myvec{x}_j\}$ & \\[-4mm]
 %row 40 
 Prover & & Verifier\\[-4mm]
 %row 41 
 Compute $\myvec{m}^{(1)}, \myvec{m}^{(1)'}, \myvec{m}^{(2)}$ as $\proto_1$,
  $\proto_1'$, and $\proto_2$ & & \\ [-7mm]
 %row 42
 Compute $C_1 \leftarrow \com(\myvec{x}, \myvec{y}, \myvec{T}, \myvec{m}^{(1)}, \myvec{m}^{(1)'}, \myvec{m}^{(2)})$  & & \\ [-7mm]
 %row 43
 & $C_1$ & \\[-7mm]
 %row 44 
 & & $r_1, r_1', r_2 \pick \F$ \\[-7mm]
 %row 45 
 & $r_1$, $r_1'$, $r_2$ & \\[-7mm]
 %row 46 
 Compute $\myvec{h}^{(1)}, \myvec{g}^{(1)}, 
 \myvec{h}^{(1)'}, \myvec{g}^{(1)'}, 
 	\myvec{h}^{(2)}, \myvec{g}^{(2)}$ & & \\[-7mm]
 %row 47 
 in $\proto_1$, $\proto_1'$, and $\proto_2$. & & \\[-7mm]
 %row 48
 Compute $C_2 \leftarrow \com(\myvec{h}^{(1)}, \myvec{g}^{(1)}, 
 	\myvec{h}^{(1)'}, \myvec{g}^{(1)'}, 
	\myvec{h}^{(2)}, \myvec{g}^{(2)})$ & & \\[-7mm]
 %row 49 
 & $C_2$ & \\[-7mm]
 %row 50 
 & $\myvec{x}, \myvec{y}, \myvec{T}, \myvec{m}^{(1)}, \myvec{m}^{(1)'}, 
 \myvec{m}^{(2)}, \myvec{h}^{(1)}, \myvec{g}^{(1)}, 
 \myvec{h}^{(1)'}, \myvec{g}^{(1)'}, 
 \myvec{h}^{(2)}, \myvec{g}^{(2)}$ & \\[-7mm]
 %row 51 
 & & Open $C_1$ and $C_2$ using last message \\[-7mm]
 %row 52 
 & & Check $f_1$, $f_1'$, and $f_2$.  \\[-7mm]
};
\draw[line width=0.50mm,shorten <=-1.5cm,shorten >=-1.5cm] (m-2-1.north west)--(m-2-3.north east);
\draw[shorten <=-1.5cm,shorten >=-1.5cm] (m-2-1.south east)--(m-2-1.south west);
\draw[shorten <=-1.5cm,shorten >=-1.5cm] (m-2-3.south east)--(m-2-3.south west);
\draw[shorten <=-1.5cm,shorten >=-1.5cm,->,thick,black] (m-5-2.south west)--(m-5-2.south east);
%\draw[-{latex[length=3mm,width=5mm]}] (m-5-2.south west)--(m-5-2.south east);
\draw[shorten <=-1.5cm,shorten >=-1.5cm,->,thick,black] (m-7-2.south east)--(m-7-2.south west);
\draw[shorten <=-1.5cm,shorten >=-1.5cm,->,thick,black] (m-10-2.south west)--(m-10-2.south east);
\draw[line width=0.50mm,shorten <=-1.5cm,shorten >=-1.5cm] (m-16-1.north west)--(m-16-3.north east);
\draw[shorten <=-1.5cm,shorten >=-1.5cm] (m-16-1.south east)--(m-16-1.south west);
\draw[shorten <=-1.5cm,shorten >=-1.5cm] (m-16-3.south east)--(m-16-3.south west);
\draw[shorten <=-1.5cm,shorten >=-1.5cm,->,thick,black] (m-19-2.south west)--(m-19-2.south east);
\draw[shorten <=-1.5cm,shorten >=-1.5cm,->,thick,black] (m-21-2.south east)--(m-21-2.south west);
\draw[shorten <=-1.5cm,shorten >=-1.5cm,->,thick,black] (m-24-2.south west)--(m-24-2.south east);
\draw[line width=0.50mm,shorten <=-1.5cm,shorten >=-1.5cm] (m-30-1.north west)--(m-30-3.north east);
\draw[shorten <=-1.5cm,shorten >=-1.5cm,->,thick,black] (m-30-1.south west)--(m-30-1.south east);
\draw[shorten <=-1.5cm,shorten >=-1.5cm,->,thick,black] (m-30-3.south west)--(m-30-3.south east);
\draw[shorten <=-1.5cm,shorten >=-1.5cm,->,thick,black] (m-32-2.south west)--(m-32-2.south east);
\draw[shorten <=-1.5cm,shorten >=-1.5cm,->,thick,black] (m-34-2.south east)--(m-34-2.south west);
\draw[shorten <=-1.5cm,shorten >=-1.5cm,->,thick,black] (m-37-2.south west)--(m-37-2.south east);
%--- prot_cm
\draw[line width=0.50mm, shorten <=-1.5cm,shorten >=-1.5cm,] (m-40-1.north west)--(m-40-3.north east);
\draw[shorten <=-1.5cm,shorten >=-1.5cm,->,thick,black] (m-40-1.south west)--(m-40-1.south east);
\draw[shorten <=-1.5cm,shorten >=-1.5cm,->,thick,black] (m-40-3.south west)--(m-40-3.south east);
\draw[shorten <=-1.5cm,shorten >=-1.5cm,->,thick,black] (m-45-2.south east)--(m-45-2.south west);
\draw[shorten <=-1.5cm,shorten >=-1.5cm,->,thick,black] (m-43-2.south west)--(m-43-2.south east);
\draw[shorten <=-1.5cm,shorten >=-1.5cm,->,thick,black] (m-49-2.south east)--(m-49-2.south west);
\draw[shorten <=-1.5cm,shorten >=-1.5cm,->,thick,black] (m-50-2.south west)--(m-50-2.south east);
\end{tikzpicture}

}
\caption{Protocols for Example \ref{examp:ir1cs}}
\label{fig:2shots}
\end{figure*}

\begin{lemma}\label{lemma:ir1cs}
Let $R$ be a decomposable relation joined by
$k$ sub-relations with special-sound $\Sigma$-protocols.
There exists a $3$-move
special-sound I-R1CS protocol for $R$. 
\end{lemma}

\begin{proof} 
We use $\proto_s(R)$ to denote the I-R1CS protocol that sequentially
composes all $k$ sub-protocols of $R$, which performs
the check of join relation $\phi$ and the R1CS checks for
all of the sub-relations, and returns the conjunction of the Boolean
results. $\proto_s(R)$ is a complete and special-sound protocol
for $R$, and it has $3k$ moves.

We use $\proto_p(R)$
to represent the $\Sigma$-protocol which is constructed by
sequentially concatenating the first messages of all 
subprotocols of $R$ for its first messages, and similarly
constructing its second and thrid messages.
$\proto_p(R)$ has $3$-moves. Its last R1CS module
performs the check of $\phi$ and the R1CS checks of all
sub-protocols.
It is apparent that for any transcript $\proto_s(R)$,
there is a one-to-one mapping $\Psi_R$ which uniquely
constructs a transcript for $\proto_p(R)$. $\Psi_R$ can 
be directly inferred from the structure information of $R$
and its sub-protocols.

It is not hard to see that $\proto_s(R)$ and $\proto_p(R)$
are equivalent, in the sense that, for an accepted transcript $t$ 
of $\proto_s(R)$, its $\Psi_R(t)$ is also accepted by $\proto_p$;
vice versa for the inverse relation $\Psi_R^{{-1}}$.
An ``invalid" transcript (which falsely proves a 
$(\myvec{x};\myvec{w}) \not\in R$
is in $R$) has negilible probability of passing $\proto_s(R)$,
and also for $\proto_p(R)$. The soundness error of the two protcols
are the same, both linearly degrading with the number of sub-protocols,
because the R1CS check modules of the two are equivalent.

When Fiat-Shamir is applied to $\proto_p(R)$,
the soundness loss {\em independent} of the number of public coin verifier 
challenges, per \cite{SpecialFiat}.
\end{proof}

We note that our I-R1CS model (even extended to more than
$3$ messages) is a strictly (restrictive) weaker 
than the original I-R1CS \cite{AIPRam}. The key difference is that, by 
the restriction of ``join" relation, our model explicitly {\em enforces that
the second and third message of sub-protocols do not have
dependence on other sub-protocols}, only the claims of sub-protocols
might be related. This ``restriction" allows the projection between
$\proto_s(R)$ and $\proto_p(R)$. This model can still capture
a wide spectran of application scenarios, e.g., all of circuits
in our ZkregPlus system satisfy the condition. 
%and the same applies
%to the RAM protocols in \cite{AIPRam}.

\subsubsection{Lookup Table}
One drawback of the protocols presented in Figure\ \ref{fig:2shots}
is that if the lookup table size is much larger than the
statement vector, it is concretely very inefficient. We provide
two solutions for the problem. 

One solution, when the total
prover work is small (e.g., very small number of folding steps
or in the case of no folding), one can use commit-and-prove approach
to take out the commitment of all elements that need to perform
lookups, and then outside of the circuit, run an independent
zero knolwedge lookup protocol. As the lookup table is fixed for
our application scenario, it can be pre-processed. We provide
a zero knowledge version of the CQ protocol \cite{cq} in Appendix 
\ref{apdx:zk_cq}, as an contribution of independent interst.
In this case, the prover only has to pay a concrete cost
of the query table size (which is indpepenent of the size of
the container table). The protocol in \ref{apdx:zk_cq}
can be replaced by the zk-cq protocol in \cite{MatrixLookup},
which has slightly smaller proof size but higher prover cost than our
protocol.

The second solution is to adopt the commit-and-fold scheme presented in 
Mangrove \cite{Mangrove}. When the total prover work (e.g., folding
$2^{17}$ circuits of an average circuit size of $2^{16}$ each,
a lookup table of $2^{28}$ can be distributed into each step, 
incurring a negligible cost of about $3\%$). We adopt the Mangrove approach
in our implementation for its convenience. 
We note that the folding lookup scheme in ProtoStar \cite{ProtoStar} yields
similar asympototic and estimated similar concrete performance.

There is one more issue to consider: for the ZkregPlus system,
there are many individual look-up tables. Range proofs are all
achieved using lookups, and there are several different
types of range proofs: ACDFA states are $23$ bits,
input alphabet is $4$-bit, and signature IDs are $16$ bits.
The $\cp$ approach needs 2 ACDFAs, one for case insensitive and
one for case sensitive. The $\bw$ and $\sde$ combined,
needs 2 more ACDFAs.  In total, these $4$ ACDFAs, have
about $160$ million transtions, to be encoded in four separate
look-up tables. We also need to encode the $\bw$ and $\sde$
discharging instructions separately.  

We put all these separate tables into one $2$-column look-up table,
with first column indicates the sub-table ID, and the second column
indicates the value. We reserve $0$ as $\texttt{null}$
sub-table ID. The logarithmic derivative approach 
\cite{Hab22,cq,ProtoStar} supports multi-column table well.

Putting all ingredients together, we now have the formal
definition of $\Sigma$-I-R1CS enriched with lookup. It is
used to specify the step-circuit in FoldPot, essentially
adopting the SuperNova \cite{SuperNova} + CycleFold \cite{CycleFold}.
This definition is adapted from \cite[Definition 1]{AIPRam}.

\begin{definition} \label{deff:ir1cs}  
A $\Sigma$-I-R1CS is a tuple $\phi$ defined
as below:
\[
(N, n, \F; \myvec{T} \in \F^{N\times 2}, 
\ell_x, \ell_{m_1}, \ell_{m_2}, \ell_{m_3} \in \nat; 
A,B,C \in \F^{n\times m}; \chi)
\]
Here $\myvec{T}$ is a pre-processed 2-column look-up table.
$\chi: \nat \rightarrow \nat$ is a function that maps
a variable index to a sub-table ID in $\myvec{T}$.
The protocol 
defines a 3-round AIP over $\F$. The prover sends the first (including
statement) and third message, and the verifier sends
the second message of uniformally sampled challenges.
Let $\myvec{x}$ be the statement, and $\big\{\myvec{m}^{(i)}\big\}_{i=1}^{3}$
be the three round of messages.
let $m$ = $\ell_x + \ell_{m_1} + \ell_{m_2} + \ell_{m_3}$.
The R1CS checker consists of three matrics $A, B, C$ of $n$ constraints.
Let $z = \myvec{x} \concat \myvec{m}^{(1)} \concat \myvec{m}^{(2)}
	\concat \myvec{m}^{(3)}$. 
The verifier accepts if $Az \circ Bz = Cz$, and
for each $i\in [m]:$ if $\chi(i) \neq \texttt{null}$ then
$(\chi(i), z_i) \in \myvec{T}$, i.e.,
$z_i$ is in sub-table $\chi(i)$.
\end{definition}

\subsection{Commit-and-Fold}
\label{cf_zksnark}
We now briefly review the commit-and-fold scheme proposed in
Mangrove \cite{Mangrove}, which we later instantiate to support
inter-circuit data passing and batch processing of non-uniform word
problems.

The base line of FoldPot is SuperNova \cite{SuperNova} and 
CycleFold \cite{CycleFold}, which solves non-uniformity of 
circuits in folding, and delivers
$O(1)$ proof. However, we still need to address the following challenges.

\begin{enumerate}
\item We need to handle large files (e.g., $2^{27}$ bytes) in
``streaming" mode. In this case, step-circuits need to pass
large amount of input/output data. We need a concretely
more efficient schemes than the general RAM approach (e.g., \cite{AIPRam}).

\item We need to discharge a collection of files in batch, and
then provide an individual (2-stage) proof for each file.
Given an aggregated word $\myvec{s}$ which is a concatenation
of $k$ segments, i.e.,
$\myvec{s} = \myvec{s}^{(1)} \concat \myvec{s}^{(2)} \concat ... \concat \myvec{s}^{(k)}$, and the individual commitments to each $\myvec{s}^{(i)}$,
we need to link these segment commitments with the $\myvec{s}$ in circuit,
and prove all of them malware free. Here the challenge is that
these word segments can have arbitrary length.
\end{enumerate}

We first present a simplified framework of Mangrove's 
commit-and-prove, in the context of committed R1CS used by Nova,
SuperNova and CycleFold, as all of our circuits are low-degree.
Recall the Nova \cite{Nova} framework: the folding prover
is given a running instance $(\myvec{U}_i; \myvec{W}_i)$ and an incoming (step)
instance $(\myvec{u}_i; \myvec{w}_i)$, where $\myvec{W}_i$ and 
$\myvec{w}_i$ are the witnessnes,
and fold them into a new instance $(\myvec{U}_{i+1}; \myvec{W}_{i+1})$. 
A folded instance is a tuple 
$(\overline{\myvec{W}}, \overline{\myvec{E}}, u, x)$,
where $u$ and $x$ are ``public" parameters of relaxed R1CS relation,
and $\overline{\myvec{W}}$ and $\overline{\myvec{E}}$ are homomorphic
commitments (e.g., Pedersen vector commitments) to the
witness and cross terms. The folding verifier in the IVC scheme
does not have to check the opening of $\overline{\myvec{W}}$ and 
$\overline{\myvec{E}}$ at each step. Instead, the 
check is {\em delayed} to the final decider proof,
and performed {\em only once} for all folded instances. 

The key observation made by Mangrove \cite{Mangrove} is that:
the committed instance can have extra commitments on subset of witnesses.
The opening of these commitments can be also delayed to the 
final decider step. Thus, if one performs a one pass aggregation (e.g.,
using Merkle tree) of these extra commitments, a commit-and-prove
scheme like LegoSNARK \cite{LegoSnark} can be achieved.  
In our context, the commitment to the subset of witness
can be used for generating the Fiat-Shamir randomness
and can be verified in-circuit.

In Mangrove, the formalization of commit-and-fold
was made using the polynomial relations, which supports high-degree.
We re-instate the result of Mangrove in the context of 2-degree R1CS
as below.
Given a $\Sigma$-I-R1CS instance $\phi$ whose
R1CS relation is $A,B,C \in \F^{n\times m}$, where
$m$ is the number of variables. Enrich $\phi$ with
$\rho$, a subset of $[0,m-1]$ sorted in ascending order.
Given the fixed size witness vector $\myvec{W}$, we
use $\rho(\myvec{W})$ to denote the projection of $\myvec{W}$ over
the indices included in $\rho$. For example, let
$\myvec{W} = (5, 30, 20, 1)$ and $\rho = (0, 2)$, then
$\rho(\myvec{W}) = (5, 20)$.
We denote the new $\Sigma$-I-R1CS instance as 
$\phi_{\rho}$.
The following defines relaxed and committed R1CS instance.
The difference from NOVA is the addition of the
projected commitments: $\overline{\rho(\myvec{W})}$.

\begin{definition} (extending \cite[Definition 12]{Nova})
\label{def:extr1cs}
A relaxed R1CS instance of $\phi_\rho$ consists
of $\myvec{A},\myvec{B},\myvec{C} \in \F^{n\times m}$ as in $\phi$,
an error vector $\myvec{E} \in \F^m$, and a scalar $u\in F$,
and input/output vector $\myvec{x} \in \F^{\ell}$.
Let $\myvec{z} = (\myvec{W}, \myvec{x}, u)$, and
$(\myvec{z}, \myvec{W^{(\rho)}})$ 
is accepted if $\myvec{A}\myvec{z} \circ \myvec{B}\myvec{z} 
= u(\myvec{C}\myvec{z}) + \myvec{E}$
and $\myvec{W^{(\rho)}} = \rho(\myvec{W})$. 
A committed R1CS instance for $\phi_\rho$ is a tuple 
$(\overline{\myvec{W}},$
$\overline{\myvec{E}}, u, \myvec{x}, \overline{\myvec{W^{(\rho)}}}$.
\end{definition}

The Nova/SuperNova folding schemes can be modified correspondingly.
When folding $\overline{\myvec{W}}$, $\overline{\myvec{W^{(\rho)}}}$ 
can be folded correspondingly (and in CycleFold, it can be fed
to a CycleFold circuit works on Grumpkin curve similarly). 
At the final decider, the verifier checks the 
folded $\myvec{W^{(\rho)}} = \rho(\myvec{W})$, 
and they are the opening to the folded commitments.
The correctness of the scheme depends on the following
lemma, which is apparent.

\begin{lemma} Let $|\rho| = v$. Assume
$\overline{\myvec{W}}$ and $\overline{\myvec{w}}$ are 
commitments to some $\myvec{W} \in \F^n$ and $\myvec{w} \in \F^n$. Similarly
$\overline{\myvec{U}}$ and $\overline{\myvec{u}}$ are 
commitments to $\myvec{U} \in \F^v$ and
$\myvec{u} \in \F^v$. 
Let $r$ be a public coin verifier challenge,
$\myvec{W'}$ be the opening to 
$\overline{\myvec{W}} + r\overline{\myvec{w}}$,
and $\myvec{U'}$ the opening to 
$\overline{\myvec{U}} + r\overline{\myvec{u}}$.
If $\myvec{U'} = \rho(\myvec{W'})$ then the probability that
$\myvec{U} \neq \rho(\myvec{W})$ or $\myvec{u} \neq \rho(\myvec{w})$ 
is negligible.
\end{lemma}

% The the knowledge extraction of the two partial commitments simply
% follows the knowledge soundness proof of the Nova/SuperNova framework.
% 
% In the following, we give a preview of how the two witness subsets
% are used.  Given a statement to prove, one can view $\myvec{W^{(1)}}$ 
% as the ``fixed" subset of witnesses that do not depend on
% the verifier's random challenge. For instance, assume we are to prove
% that a string $w$ is accepted by a DFA, its witness will include
% the sequence of states along the acceptance path. This is regarded
% as ``fixed" witness. Then, transitions can be computed from
% source/destination states and characters from $w$. They are also
% fixed.  To assert that
% a transition $t$ is valid, this can be achieved by proving that
% $t$ is contained in a preprocessed lookup table. We apply
% the logarithmic derivative technique \cite{Hab22} as shown in $\proto_1$ of
% Figure\ \ref{fig:2shots}. This needs a verifier challenge. Applying
% Fiat-Shamir, the challenge is generated in-circuit by hash on
% the prior messages (which is the set of ``fixed" witness). Then,
% more witnesses are included, but they depend on the Fiat-Shaimir 
% randomness, which in turn depends on the set of fixed witness.
% Later, we will see that for supporting the batch discharging of
% file collections, we need two layers of Fiat-Shamir. Then 
% $\myvec{W^{(2)}}$ depends on $\myvec{W^{(1)}}$, and the
% rest of witness depend on $\myvec{W^{(2)}}$.
To apply the finger-printing technique in circuit,
a binding commitment to the fixed witness subsets have to be
generated by the prover by performing a 1st walk before generating
the proof, so that the commitment can be used to generate Fiat-Shamir
randomness. Following Mangrove, but in a sequential folding scheme,
we use a hash-chain.

Let their be $k$ step-circuits involved in a prover job.
Let the $\myvec{W^{(\rho)}} = 
	\myvec{W^{(\rho,1)}} \concat
	\myvec{W^{(\rho,2)}} \concat ...
	\myvec{W^{(\rho,k)}}$.
We define $\hashchain(\myvec{W^{(\rho)}}, i)$ recursively as:
\[
	\left\{
		\begin{array}{l l}
			\hash(\overline{\myvec{W}^{(\rho,1)}}) & \textrm{ if } i=1 \\
			\hash( \hashchain(\myvec{W^{(\rho)}}, i-1), 
				\overline{\myvec{W}^{(\rho,i)}})  & \textrm{ if } i>1 \\
		\end{array}
	\right.
\]
Now that the hashchain can be computed in circuit by slightly
enhancing the Nova framework because
at each step $\overline{\myvec{W}^{(1,i)}}$ is available
for random linear combination during folding (and also
as input to the CycleFold component on Grumpkin).


\subsection{Passing Data Between Circuits}
The first application of 
$\hashchain(\overline{\myvec{W}^{\rho}})$ is to pass data
between circuits. Recall that $\myvec{W}^{\rho}$ serves
as a ``fixed" memory.  Here we require that input and output
buffers of circuits are located at fixed positions of
$\myvec{W}^{\rho}$ and the hashchain serves as the seed
for generating Fiat-Shamir randoms. We first handle
the problem of passing data between $k$ uniform circuits.

For each circuit $i$, its its fixed chunk 
$\myvec{W}^{(\rho,i)}$ has a fixed
layout of $u$ elements allocated for input and also $u$ elements
for output. We denote these two regions as
$\myvec{I}^{(i)}$ and $\myvec{O}^{(i)}$.  
Note that $r = \hashchain(\myvec{W}^{(\rho)})$.

We then let $k$ circuits jointly perform the following computation
(in circuit), essentially asserting the equivlance of
vector 
$(\myvec{I}^{(1)}, ..., \myvec{I}^{(k-1)})$ and
$(\myvec{O}^{(2)}, ..., \myvec{I}^{(k)})$
\[
  \sum_{i=1}^{k-1} \sum_{j=1}^u r^{(i-1)*k+j} \myvec{I}^{(i)}_j = 
  \sum_{i=2}^k \sum_{j=1}^u r^{(i-2)*k+j} \myvec{O}^{(i)}_j 
\]

The above can be computed by accuulating partial results 
across circuits, and the equality check be performed in the
last circuit. The public coin verifier challenge $r$
guarantees that with negligible probability
the input and output buffers do not match.

We next extend the above idea to non-uniform circuits in SuperNova. 
Each circuit (type) has its own configuration (including the
size of pre-allocated buffer for its input/output regions). These
are {\em fixed constants} embedded in the circuits.
Then we need to keep track of the actual
input/output size. Let them be
$u_{in,i}$ and $u_{out,i}$ for circuit $i$, they are
located at some fixed positions in In $\myvec{W}^{(\rho,i)}$.
Then we just need to perform the following check using
all circuits in folding:
\[
  \sum_{i=1}^{k-1} \sum_{j=1}^{u_{in,i}} r^{(i-1)*k+j} \myvec{I}^{(i)}_j = 
  \sum_{i=2}^k \sum_{j=1}^{u_{out,i}} r^{(i-2)*k+j} \myvec{O}^{(i)}_j 
\]
In addition, each circuit needs to perform in-circuit for: (1) 
$u_{in,i}$ does not exceed to static input buffer limit; and (2)
total size of read from the input buffer does not exceed $u_{in,i}$.
Then similar checks apply to $u_{out,i}$ for each circuit.
This approach fits our specific application setting, and
has better concrete performance (mainly $1\times$ the input
and output buffer size), compared with the
the general RAM approach (e.g., \cite{AIPRam}, which
costs at least $43\times$ of memory cell accesses). 

\subsection{Batch Processing with Non-uniform Circuits}
We now discuss how to batch process a large collection of non-uniform
files via folding. By ``non-uniform" we mean that these files have
different sizes, and may generate witness of different size, thus
relying on non-uniform circuits. In practice, we define
one or more circuits for each integer logarithm of file sizes.
We first formally define the problem.
Let $\myvec{s} = \myvec{s}^{(0)} \concat ... \concat \myvec{s}^{(k)}$
be a concatenation of $k$ words $\myvec{s}^{(i)}$, and let 
$\myvec{\ell} = \left(|\myvec{s}^{(i)}|\right)_{i=1}^k$ be the vector
of string lengths.
We call $\myvec{s}^{(i)}$
the $i$'th {\em string fragment} and $s$ the {\em batched string}.

Equation\ \ref{eqn:batch} presents the relation,
where
$\left(k, \overline{\myvec{s}}, \left(\overline{\myvec{s}^{(1)}}, ..., \overline{\myvec{s}^{(k)}}\right)_{i=1}^k, \overline{\myvec{\ell}}, R\right)$
are the public statement also visible to the verifier but 
$\left(\myvec{s}, \myvec{\ell}, \left(\myvec{s}^{(i)}\right)_{i=1}^k\right)$
are the secret witness, which are separated from the public
statement using ``\texttt{;}". 
Here $R$ represents a collection of regex signatures.
\begin{equation}
\left\{
\begin{array}{l}
	\left(k, \overline{\myvec{s}}, 
		\left(\overline{\myvec{s}^{(1)}}, ..., \overline{\myvec{s}^{(k)}}\right)_{i=1}^k,
		\overline{\myvec{\ell}}, R; \myvec{s}, \myvec{\ell},
		\left(\myvec{s}^{(i)}\right)_{i=1}^k
	\right): \\
	(1) \forall i\in [k]: \myvec{\ell}_i = |\myvec{s^{(i)}}| \And \\
	(2) \myvec{s} = \myvec{s}^{(1)} \concat ... \concat \myvec{s}^{(k)} \And \\
	(3) \forall i\in [k]: \forall r\in R: \myvec{s}^{(i)} \not\in L(r)
\end{array}
\right\}
\label{eqn:batch}
\end{equation}

Intuitively, the public statement asserts the following
given the commitments to all string fragments and length information:
(1) All strings $\myvec{s}^{(i)}$ 
are malware free\footnote{Note that this is however, different from
asserting that $\myvec{s}$ as a single string is malware free.
When $\myvec{s}$ is a concatenation of many files, it is easy
to trigger false positive (identified as a malware).},
and (2) Each $\overline{\myvec{s}^{(i)}}$ indeed hides the $i$'th
string fragment of $\myvec{s}$ as defined by the length vector,
which hides inside $\overline{\myvec{\ell}}$.
Naturally, the proof consists of two parts: a global batch proof
for all, and an individual proof for each string fragment.
Our solution provides $O(1)$ proof size for both.

Note that unlike the folding scheme where Pedersen Vector Commitment
is used, here in the problem statement
all commitments are {\em KZG commitments} 
(where a vector of field elements serve
as the co-efficients in the univariate polynomial for KZG).
Given a KZG commitment $\overline{\myvec{x}}$ 
to a vector $\myvec{x}$
of degree $d$, we 
use $\prfkzg(\overline{\myvec{x}}, a, y)$ 
as the KZG polynomial
opening proof for: $y = \sum_{i=0}^d \myvec{x}_i a^i$.
It is known that given the Pedersen Vector Commitment and
KZG commitment to the same vector, the Quasi-linear NIZK (QA-NIZK)
\cite{Gonzalez15} provides linear prover complexity and constant
size proof for their equivalence, as demonstrated in LegoSnark \cite{LegoSnark}.

We now address (1) and (2) of the relation and delay (3) to the next
subsection.  Due to the
use of non-uniform circuits, we discharge {\em at most} one 
string fragment each circuit.\footnote{In another word,
no circuit will process multiple string fragments.}
For large string fragment, it may span over multiple circuits.
For $\myvec{s}^{(i)}$ spanning over $d$ circuits, we list
these circuits as $\circuit^{{i,1}}, ... \circuit^{{i,d}}$.
Similarly, their fixed witness subset is denoted as
$\fm{i,1}, ..., \fm{i,d}$.

The prover first walks over the public statement and defines
the $\fm{i,j}$ for each step-circuit $\circuit^{(i,j)}$.
Note that this is ``fixed" information that does not depend
on any verifier challenge. It can be directly generated from
the problem statement.
\[
	\fm{i,j} = \left(\myvec{s}^{(i,j)}, \myvec{\ell}^{(i,j)}, 
	\myvec{r}_{i}, \myvec{v}^{(i,j)}\right)
\]
Here $\myvec{s}^{(i,j)}$ is the $j$'th slice of $\myvec{s}^{(i)}$
allocated to circuit $\circuit^{(i,j)}$.
$\ell^{(i,j)} = \sum_{k=1}^j \myvec{s}^{(i,k)}$ is the accumulated
length of slices up to $j$. Let $\circuit^{(i,d)}$ be the
last circuit for processing $\myvec{s}^{(i)}$, then
$\myvec{\ell}_{i,d} = \ell_{i}$, which is checked in circuit.
$\myvec{r}_i$ is a Fiat-Shamir verifier challenge used in
the individual proof. It is pre-computed 
outside circuit, and served as a non-deterministic advice:
\begin{equation}
	\myvec{r}_i = 
		\rohash\left(\overline{\myvec{s}}, i,
			\overline{\myvec{s}^{(i)}}, 
			\overline{\myvec{\ell}}\right) 
\end{equation}

Define $\myvec{v}_i$ as the following, where
$\myvec{s}^{(i)}_u$ is the $u$'th element of
string fragment $\myvec{s}^{(i)}$:
\begin{equation}
	\myvec{v}_i = \sum_{u=0}^{d} \myvec{r}_i^u \myvec{s}^{(i)}_u
\end{equation}

Apparently, $\myvec{v}_i$ is an evaluation of $\myvec{s}^{(i)}$
as the co-efficient vector of a polynomial.
$\myvec{v}_i$ can be be computed by accumulating its value
across circuits. This is done by computing
$\myvec{v}^{(i,j)}$ as the following:
$\sum_{u=0}^{\ell^{i,j}} {\myvec{r}}_i^u {\myvec{s}}^{(i)}_u$,
and apparently for the last circuit in the sequence,
$\myvec{v}^{(i,d)} = \myvec{v}^{(i)}$.

 The prover then computes the hashchain of the fixed witness subset,
 and based on it and other KZG commitments in the public statement,
 compute the Fiat-Shamir challenge $c$:
 \[
 	c = \rohash\left(
 		\hashchain\left(\emptyfm\right), 
 			\overline{\myvec{r}}, \overline{\myvec{v}}
 	\right)
 \]

Then the prover computes the rest of the witness fo all circuits,
based on challenge $c$, and after all are folded, the last circuit
outputs the following:
\begin{enumerate}
	\item $y_v = \sum_{i=1}^k \myvec{v}_i c^i$
	\item $y_r = \sum_{i=1}^k \myvec{r}_i c^i$
	\item $y_s = \sum_{i=1}^n \myvec{s}_i c^i$
	\item $y_\ell = \sum_{i=1}^k \myvec{\ell}_i c^i$
\end{enumerate}

Each individual circuits performs the verification and asserts that
$\myvec{v}^{(i,j)}$ are computed inside the circuit correctly
using $\myvec{r}_i$ and $\myvec{s}^{(i,j)}$, and similarly
$\myvec{\ell}_i$ accumulated correctly using $\myvec{\ell}^{(i,j)}$.
Then the global {\bf batch proof} consists of the following.
Thus, by asserting the values of
the in-circuit $y_v, y_r, y_s, y_\ell$ and their counter-parts
in the KZG opening proof, the verifier is convinced
that the shares of $\myvec{v}, \myvec{r}, \myvec{s}, \myvec{\ell}$
are distributed across all circuits in their fixed memory correctly.

\begin{enumerate}
\item $\prf_{\texttt{snark}}$ is the final decider proof (compressed
	using Groth16 \cite{Groth16} in CycleFold) for the folded circuits.
\item 
$\prfkzg(\overline{\myvec{r}}, c, y_c)$,
$\prfkzg(\overline{\myvec{s}}, c, y_s)$
$\prfkzg(\overline{\myvec{v}}, c, y_v)$, and
$\prfkzg(\overline{\myvec{\ell}}, c, y_\ell)$
which can be further aggregated.
\end{enumerate}


We now proceed to each {\bf individual proof}. 
Given a homomoric commitment $\overline{\myvec{x}}$ to a vector,
$\prfvec(\overline{\myvec{x}},i,y)$ asserts that the $i$'th
element of $\myvec{x}$ is $y$. Then for each
statement $\left(\overline{\myvec{s}^{(i)}}, i, \overline{\myvec{s}}, 
\overline{\myvec{\ell}}\right)$, i.e.,
$\overline{\myvec{s}^{(i)}}$ hides the $i$'th string segment in $\myvec{s}$
as defined by the length vector in $\overline{\myvec{\ell}}$, 
the individual proof consists of the following (in addition
to the global batch proof):
\begin{enumerate}
	\item Fiat-shamir transform of verifier challenge: $\myvec{r}_i = 
		\rohash\left(
		\overline{\myvec{s}}, \overline{\myvec{\ell}},
		\overline{\myvec{s^{(i)}}}, \overline{\myvec{v}}\right)$.
	\item 
		$\prfvec(\overline{\myvec{v}}, i, \myvec{v}_i)$,
		$\prfvec(\overline{\myvec{r}}, i, \myvec{r}_i)$
	\item $\prfkzg(\overline{\myvec{s}^{(i)}}, \myvec{r}_i, \myvec{v}_i)$,
\end{enumerate}

Intuitively, the proof works by
 performing finger-printing
 on $\myvec{s}^{(i)}$ at random point $\myvec{r}_i$ (producing
 $\myvec{v}_i$).\footnote{In practice, the KZG commitment
 of $\myvec{s}^{(i)}$ by padding it with extra random elements,
 as many systems such as Halo2 \cite{Halo2}.}
%  Note that actually $\hashchain(\emptyfm)$ is {\em not} needed
%  to generate $\myvec{r}_i$, because the contents $\myvec{s}$
%  have already been ``fixed" and linked to be equivalent to
%  $\overline{\myvec{s}}$ in the global batch proof, which eliminates
%  the possibility of adversary action to adaptively 
%  set the conntents of $\emptyfm$ (including
%  $\myvec{s}^{(i)})$ based on $\myvec{r}_i$.
%  \footnote{If $\hashchain(\fm)$
%  is used to generate $\myvec{r}_i$, then a second layer
%  of $\emptyfm$ is needed, by increasing one CycleFold gadget for each
%  circuit. At the final decider step, this incurs a cost of
%  $5$ million R1CS per non-uniform circuit type.}
 
\subsection{==== NOT USED BELOW====}



implementation of folding schemes, 

In this section, we discuss how to efficiently encode the
5-stage regex discharging algorithms in zkSNARK. Given that
there are many small files (e.g., emails), it is not economical
to provide a proof with size almost comparable to the original
file size. We opt for pairing based constant size proof system
such as Groth16 (instead of log-size proof system such as sumcheck
based systems). Our choice of proving scheme s is also limited
by the choice of open-source folding scheme implementations available
at the time the project is implemented.
In this case, we choose to adopt Sonobe (an implementation of
Nova + Cyclefold folding schemes that work for BN254-Grumpkin curve
that allows producing final decider proof in Groth16). 
,It is possible to encode our 5-stage discharging algorithms using
many other schemes, such as Plonk based for universal prover keys,
or Mangrove for PCD foldings schemes with constant memory footprints,
and we leave them for future work.

We have several design challenges to overcome.
First of all, we would like to take advantage of
the folding schemes so that the prover cost can be
minimized. Second, in addition to mass processing regex signatures,
we can also mass-discharge files, e.g., providing
a commitment to a large collection of files (e.g., all binary executables
or all Enron emails), provide an $O(1)$ proof for free of
malware for all of them, and then provide a membership 
of an individual file commitment belong to the combined commitment.
Third, to achieve the aforementioned scheme, we have to
handle vast difference in file sizes from $2^{9}$ to $2^{27}$.
This concerns handling non-uniform circuits (with different size).
Fourthly, we need to handle lookup arguments effectively.
Although CQ lookup is handled in ProtoStar, we present
a simpler stand-alone zk-CQ protocol which shows similar
concrete prover effiicency.

\subsection{Consideration of Folding Cost}
For the vast majority of folding schemes that follow the Nova folding,
it is shown that prover concrete efficiency is $4\times$ of Groth16.
However, there is a final decider scheme (in Sonobe) that is expensive
to compute, its cost is shown in the following, given $n$ the step-wise
circuit size, $k$ the number of steps applied in folding: 
\[
	2^{13.5} + 3n
\]
Mainly, the $2^{13.5}$ is a fixed non-native field arithmetic component in the decider and $3n$ is the linear cost on proving the step-wise R1CS relation.
Now given that folding itself takes $2n$ group operations per step, the fraction of the decider proof genreation among the total cost is, when $k$ is large enough:
\[
\frac{2^{13.5} + 3n}{2nk} \approx \frac{1.5}{k}
\]

In our later presentation, we will need to customize the Sonobe final decider
proof, which will add several more ``expensive" constant size computation. 
However, given that $k$ is large enough, they will be amortized soon, which
does not cause concern. 

Also another interesting point is the step-wise circuit size.
In the celebrative work of Mangrove, it is shown how any arbitrary Plonk
Irelation can be chunked into any arbitrary thread of smaller computation,
which allows arbitrary folds of scalability via parallelism, i.e.,
it allows chunking a large R1CS into many very small steps.
In our case, server RAM is not a concern, so that we allow large
step-wise circuit size, as long as it does not make $k$ too small.

\subsection{Road Map}
In this section, we will first discuss how to encode the $5$-stage 
regex discharging algorithm efficiently in R1CS. We then discuss a 
customized folding scheme to support batch process a large collection
{files. 

\subsection{Encoding $5$-Stage Regex Discharging Algorithm}
\subsubsection{Encode DFA (TODO)}
In practice, an ACDFA can be transformed into a DFA. As DFA acceptance is
needed in various parts of the entire algorithm, we discuss them first.

The key idea is that: there is an external and precomputed lookup table
of all DFAs needed (notice that there are multiple tables). But we integrate
them into one table of two columns $\texttt{TableID}, \texttt{Entry}$.

TODOS in details: rough idea

(1) discuss how to encode DFA transition
(
(2) discuss what kind of information will be stored in external table, e.g., transitions (for different DFAs), the critical patterns needed by each signature).

(3) present an algorithm in R1CS or circuit, which takes in witnesses: (source state, char, dest state) sequence, produces (still secret) output: the set of signatures that should be triggered (those who have their critical patterns appear along the acceptance path)

(4) present a variety of gadgets that proves permutation of vectors, equivalence of vectors and its sorted form, packing a vector to a set etc. In particular, notice that we can use the in-circuit form of the Habock scheme \cite{Hab22} for proving lookup relation between two vectors (when they are small) {\em in circuit}, this is useful for proving a ``set" is reduced from a vector with duplicate items.

(5) Given the gadgets, this motivates the definition of R1CS relation with $2$ stages (to allow to generate the random needed by the gadget), and 
the need to tag some some elements for the use of external lookup argument.

\subsection{Customized zkSNARK with Folding Schemes}
To support the design goal of batch processing large quantity of files against
large collection of regex signatures, we present a customized
folding scheme by adapting the Nova \cite{Nova} + Cyclefold \cite{CycleFold}
with a Groth16 final decider proof generator (based over the Sonobe
\cite{Sonobe} framework).

\subsubsection{Discussion of Fiat-Shamir}
Applying Fiat-Shamir is tricky, as pointed in \cite{WeakFS},
Fiat-Shamir needs to include statement to defeat adaptive adversaries.
There were other attacks on parallel repetition of multi-round
protocols, however, it is shown in \cite{SpecialFiat} that
for {\em special sound} interactive protocols the security degrades
only linearly with the number of queries to the random oracle,
thus, {\em independent of} the round of an interactive protocol.\footnote{The intuition is that since the interactive protocol is special sound (allowing
efficient knowledge extractor given a transcript tree to extract true
witness). If the attacker of Fiat-Shamir has non-negligible probability
of success, then the attacker is able to inject too many fake 
but accepted transcripts, thus making the true knowledge extractor
unable to extract true witness efficiently.}
This fomrs the basis of security of application of Fiat-Shamir
in folding schemes such as ProtoStar \cite{ProtoStar}
and Mangrove \cite{Mangrove}.

\subsubsection{Discussion of R1CS Model}
We assume the following extended R1CS model: the variables each carry
a (an optional) collection of tags which indicates the range proofs 
it need to satisfy and up to one CQ lookup. The CQ lookup table
is a 2-column table with the first column indiates \texttt{table\_id},
the the 2nd conlumn containing the entry values. Thus, we are
allowed to concatenate multiple tables into one 2-column CQ table.
In particular, we have 4 ACDFA's transition tables (for $\cp$
and $\bw$ and each needs two for supporting ignore-case). We also
need to encode the $\bw$ entries for all signatures and $\sde$
entries, thus the total CQ lookup table is close to $2^{28}$ (which is
the maximum size allowed by the BN254-Grumpkin curve pair, because
it is the largest factor of $|F_r|-1$, needed for FFT).
Note that this also restricts circuit size to $2^{28}$. (Question:
can we find a similar curve pair like BN254-Grumpkin for
BLS12-381?)
{
\subsubsection{Discussion on Lookups}
When the lookup table is huge, it is not a good idea to embed
it in the circuit (because it is part of the circuit size).
In ProtoStar \cite{ProtoStar}, a folding scheme is supported for
folding the Hab\"{o}ck scheme \cite{Hab22}. The prover pays
a cost indpendent of the lookup table size per step of folding,
however, this cost is eventually paid in the construction of final
decider proof. For our scheme, the total cost paid by the prover is
list below, assuming query table size is $m$ and lookup table size
is $N$, for $u$ steps of computation:
\[
	u\times 3m + |N|
\]
If one uses pre-computed quotient cache like \cite{cq}, this can
be improved to (in terms of group operations):
\[
	u\times 3m + 8\times \min(u\times m, |T|)
\]
Here for each folding operation, because there is a need to
encode three vectors $\vec{h}$, $\vec{g}$ and (a supppresed form of)
of $\vec{m}$ (because $\vec{m}$ is sparse) in the transcripts,
when calculating their Pedersen commitment, the prover has to pay
$3m$ group operations (in addition to the other witness variables).\footnote{Here we ignore the speed-up caused by MSMs for simplicity of discussion.}

In the final proof generation, the checks of three equations
as listed in the last step of section 4.3 of ProtoStar has to
be represented in R1CS. Given that now the full entries of 
$\vec{m}$ has to be reasoned about, its size is $\min(um, |T|)$.
In our setting, we generate Groth16 proof, this is equivalent
to $8$ times of G1 operations (which is 3 G1 and 1 G2 for each R1CS
constraint, where 1 G2 costs 4 times of G1 operation).
When $um < |T|$, this amounts to $11\times um$ group operations
in total for the prover.

We later introduce a concretely more efficient scheme than the above.
We construct a complete zk protocol for the CQ protocol. Its prover cost is
$8m$ G1 operations for each step of computation, thus totaling
$8um$. We show that its folding cost is negligible.  We commit to each
smaller query table at each step, this is a part of the public input.
The verifier circuit just performs the random linear compbination
of this group element (which costs one group operation over 
Grumpkin curve, which is about $1000$ multiplications). Then in the final
decision step, a fixed 10M R1CS is paid for the verification of the
group operation, and a fixed $m$ R1CS constraints is paid to verify
the fragment of witness corresponds to the query table commitment.
In total we pay:
\[
	8um + 8\times(2^{23.5} + m)
\]
It is not hard to see that when $u$ is reasonably large, it pays off
for the constant cost in decider quickly. For instance, set $m = 2^{20}$,
when $u>20$, it already makes even for $3um>8\times(2^{23.5} + m)$.
When $u=1000$, our scheme's total cost is about $73\%$ of the ProtoStar
scheme.

In Mangrove \cite{Mangrove}, the cost of $|T|$ is directly
distributed into the chunk computation. It might work in our context where
the chunk size is big and the number of steps of big, thus the
amortized cost compared with step-wise circuit being folded
is small. Its total cost is:
\[
	4um + 4T/u*u + (8u + 8T/u)*u + 8(T/u + 3m)
	\approx 4um + 8u^2 + 12T
\]
Details: when folding, the extra witness length has to cover
$\vec{h}$, $\vec{g}$, $\vec{m}$ and the part of table $T$ entries
distributed (this is because in its secret witness it has to allocate
the chunk of corresponding $T$ entries to avoid switch-case logic
on big $T$ contents itself in the circuit). Thus they incur extra length: 
$m + m + T/u + T/u$, and given NOVA cost of 2MSM,
it is: $4m + 4T/u$. 

In the recursive (verifier) circuit, given public input $i$,
and given the commitment to the $T$ chunk, the circuit has to
use a lookup (or membership) 
argument itself to verify that this is correct (then
the folded argument will be verified in the final decision proof).
So this needs an O(u) circuit, which translates to 8u group operations.
Also to compute the commitment, it needs 8T/u.

For the final decider proof, it needs to verify (1 time)
the folded relation for: (1) computing of the 
commitment of combined $T$ chunk of $T/u$ size; (2)
combined relation to compute $\vec{h}$, $\vec{v}$ and $\vec{m}$ (chunked),
however, unlike the ProtoStar and our case, it's size $m$.
In total, this incurs one time cost of $8\times (T/u + 3m)$.
In summary: when $um<T$, our zk-cq approach is better.
When $T<um$ and $u<m$, the Mangrove approach is better.

\subsection{Dynamic and Small In-Circuit Lookups}
There are various situations that we need to reason about the range
of witness elements, e.g., range of source/destination states and
char width (4-bits). They can be processed conveniently with low cost
using lookup arguments. There are also cases that
``dynamic" lookup arguments are needed when the lookup table itself
is not fixed, e.g., for proving one vector is a subset
of another. In this case, there are many choices of encoding available
for in-circuit lookups: plookup \cite{plookup} and its variation
Halo2 lookup \cite{Halo2}, Hab\"{0}ck's scheme \cite{Hab22}.
In terms of ``extra" witness variables they introduced (thus
impacting the folding cost of 2MSM, which is based on the
number of witness variables), letting $m$ and $n$ be the
query and lookup table size, the cost are ({\bf TO VERIFY LATER}):
$3(m+n)$,  $3m+n$, $2m+n$, respectively. In this case, we
implement in-circuit lookup using Hab\"{o}ck scheme \cite{Hab22}.
For In-Circuit lookups, they are encoded as part of the step-wise
R1CS, thus, does not cause any much extra considering decider cost
is amortized by steps.

There is still one problem regarding: small range proofs. Shall we
do in-circuit lookups or cq? Given that both small (dynamic) lookup
and externdal CQ lookup both exists. If we do range proofs via
in-circuit lookup, it addes the folding cost: $2m+n$ (Hab\"{o}ck) per
step, thus totaling: $u \times 2 \times (2m+n) \approx 6um$,
For CQ lookup, it actually adds ``nothing", because we do not add
any additional witness, however, we do have to pay the final cost
for all of them: $8um$, but given that it's actually 2 column table,
we are actually paying $16um$, the in-circuit wins (for range proof).

\subsection{Implementation Consideration}
There are several choices of encoding schemes, e.g.,
Circom \cite{circom}. However, considering that
we are instrumenting the standard NOVA and R1CS model,
using Arkworks library \cite{arkworks} and instrumenting the Sonobe
library seems a more convenient choice for us.

\subsection{Supporting Batch File Discharging}
Consider the scenario that there are many files of the same
size category (e.g., over $100k$ emails of size $2^{16}$).
We would to first leverage folding to speed up prover cost,
and generate one proof to prove that all of them are free of
malware. Also, for each file commitment, we provide a membership
proof which shows that it belongs to the commitment to the
entire file collection. Thus, for an individual file, its proof
consists of 2 parts, and we would like both to be $O(1)$.

We first need a scheme for ``commitment of commitments". 
There are several alternatives available, e.g., 
AFGHO commitment used in generalized inner product argument \cite{Bunz21},
which has log size proof.  Another potential solution is
the bivariate poly commitment used in Pianist \cite{Pianist},
however, we do not see a solution provides membership proof
better than $O(mn)$ (where $m$ is the number of words and $n$ is the
max length of each word).

We use the solution below.  The prover precomputes all $[f_i(s)]_1$
as Pedersen vector commitment $\sum_{i=1}^n f_{i_j} [\lambda_i(s)]_1$.
Each such commitment can be written as $(x_i, y_i)$ over $\F_q$. 
Then the prover can compute similarly commitment to $\vec{x}$
and $\vec{y}$, let them be $\com(\vec{x})$ and $\com(\vec{y})$.
It is known that to ``mass-produce" the membership for
all $x_i$ belonging to $\com(\vec{x})$ is $O(m\log(m))$ (basically
the batch KZG proof generation used by CQ \cite{cq} using
\cite{FK23}). So, each membership proof is $O(1)$, and the
amortized cost for each proof is $O(\log(m))$. Note that the commitment
has to be performed over BN254. If it's on Grumpkin, since
there is no pairing available, it needs inner product argument (IPA)
like argument, and this would need log size proof.

We then need to argue that at each step $i$, it is indeed using 
$(x_i, y_i)$, a point in G1.
At each step, the step-circuit can compute $\lambda_i(r)$ based
on $i$ and $r$ and the Barcentric Langurange poly value $l(r)$.
The weight factor $w_i$ of Barcentric can be stored as a small lookup table.
Then each step-circuit can compute $\lambda_i(r)$ in $O(1)$ time.

Each step circuit has $x_i$ as secret witness input. It applies and
accumulates $x_i \lambda_i(r)$ (passing the partial sum as $z_i$
to next circuit). So the final output contains 
$\sum_{i=1}^m x_i \lambda_i(r)$. Similarly can be done for the
$y_i$ series. Then the verifier can check the value of
$\sum_{i=1}^m x_i \lambda_i(r)$ via a KZG proof. 
This proves that the $(x_i, y_i)$ at each step
is indeed the ones contained in the entire commitment
$\com(\vec{x})$ and $\com(\vec{y})$.

There is one problem though, the field arithmetic of $(x_i, y_i)$
is non-native. So like Sonobe for group element operation,
this is ported to the Grumpkin curve: it has a separate
R1CS relation, which takes: $x_i$, $y_i$, $w_i$, $l(r)$ 
and performs the calculation, but the look-up of $w_i$ is performed
over BN254. This circuit is much smaller than the group element
exponentiations and only costs $<20$ R1CS constraints.

Then there is a one time in the decision proof which similarly
verifies the non-native R1CS relation over BN254. But much smaller
than the group operation cost.

{\bf TODO: give concrete cost.}.

\subsection{zk-CQ and Integration}
Use zk-CQ only for SINGLE instance (with no folding) lookup argument.
And this saves most of the cost and has a good justification, and
there is no other solution good than it for small circuit size.

For foldig, simply use Mangrove approach. Here, as mentioned earlier,
the segment of table $T$ needs to be stored separately. Then,
we need to verify the segment commitment is commensurate with that
part of the secret input. This membership proof can be done similarlty
to the argument of word membership belong to a large collection.
Then we are accumulating as step input and output.

\subsection{Handle State In and Output}
For very large files, we have to split into multiple steps, especially
for side. Then there are state information to pass between steps,
however, we want to save the hash operation.

Let's say, we reserve a special segment out of witness as a state output
and also another segment as output.
Then each step circuit outputs a Commitment to that segment.
The next step takes it in, and verifies that its segment of input
matches that of the preivous output. This allows to pass LARGE state,
and the verification relation can be folded and is ONLY needed
to be performed once in the final decision proof. Also the
commitment is not restricted, as long as it is homomorphic, e.g.,
KZG or Pedersen vector commitment, or maybe even coding based
or lattice based as shown in LatticeFold.

More formally, let $\vec{w}$ be a vector of witness and let
$\sigma$ be map so that $\sigma(\vec{w})$ is a fixed sub-vector
of $\vec{w}$ given a fixed list of indices (e.g., given $\vec{w}$
100 elements, $\sigma$ retrieves the fixed list of $1$st and $50$th
elements from it and forms a 2-element sub-vector.)
Given the same $\sigma$, and let $\vec{w}$ and $\vec{v}$ be
two vectors of the same length (or even different length
but allowing valid $\sigma$). It is not hard to see that,
given a verifier sampled challenge $r$,
$\com(\vec{w}) + r\com(\vec{v})$ encodes a vector
$\vec{w} + r\vec{v}$ and 
$\com(\sigma(\vec{w})) + r\com(\sigma(\vec{v}))$
encodes $\sigma(\vec{w} + r\vec{v})$.

Using the homomorphic property, in the committed R1CS relation,
we also disclose $\com(\sigma(\vec{w}))$ (with 
$\com(\vec{w})$. In the final decision, we just need to perform
the check of $\sigma$ once. And there is only one additional
group operation in recursion overhead. Using this approach,
we can pass large state variables between states.
{\bf TO CONFIRM:} this seems MUCH less expensive than
maintaining a RAM memory as in zk-EVM machines, because
we do not maintain the RAM for evyer single EVM instruction,
but for large chunk of step-wise R1CS.

\subsection{Summary on Folding Scheme}
In summary we need to modify the Sonobe framework for:
(1) additional elements for supporting word colleciton
(2) to support state variables.

The development work will: 
(1) support individual file verification
(2) fold large collections of small emails
(3) support state variables in large circuits, and non-uniform circuits.



